\documentclass[12pt]{article}

% ---------- packages ----------
\usepackage[utf8]{inputenc}
\usepackage[T1]{fontenc}
\usepackage{amsmath,amssymb}
\usepackage{graphicx}
\usepackage[margin=1in]{geometry}
\usepackage{microtype}

% ---------- title ----------
\title{\textbf{Physics-Informed Deep Learning for Noise-Robust IVIM Parameter Estimation in Breast DW-MRI}}
\author{}
\date{}

\begin{document}
\maketitle

\section{Introduction}

Breast cancer remains the most prevalent malignancy among women worldwide, accounting for approximately 2.3 million new diagnoses annually, with early and accurate detection being critical for improving patient outcomes~\cite{sung2021}. Diffusion-weighted magnetic resonance imaging (DW-MRI) has emerged as a powerful non-invasive tool for characterizing breast tissue by probing the microscopic motion of water molecules within biological tissue~\cite{lebihan1986}. Unlike conventional anatomical imaging, DW-MRI provides quantitative physiological information that can distinguish between benign and malignant lesions without the need for contrast agents, making it particularly valuable in clinical settings where gadolinium-based contrast may be contraindicated.

The Intravoxel Incoherent Motion (IVIM) model, originally proposed by Le Bihan et al.~\cite{lebihan1988}, extends the standard diffusion framework by separating the observed signal attenuation into two distinct components: true molecular diffusion of water in tissue and pseudo-diffusion arising from microcirculation of blood in the capillary network. The IVIM signal is described by a biexponential decay model:
%
\begin{equation}
\frac{S(b)}{S(0)} = (1 - f) \cdot \exp(-b \cdot D_t) + f \cdot \exp(-b \cdot D^{*})
\label{eq:ivim}
\end{equation}

where $S(b)$ is the measured DW-MRI signal at diffusion weighting factor $b$ (s/mm\textsuperscript{2}), $S(0)$ is the baseline signal without diffusion weighting, $f$ is the perfusion fraction representing the volume of blood flowing in the capillary network relative to total water content, $D_t$ is the true tissue diffusion coefficient (mm\textsuperscript{2}/s), and $D^{*}$ is the pseudo-diffusion coefficient (mm\textsuperscript{2}/s) reflecting the incoherent microcirculation of blood.

These three biophysical parameters carry distinct clinical significance for breast cancer characterization. The perfusion fraction $f$ is elevated in malignant tumors due to tumor-induced angiogenesis, which increases the density of the microvasculature~\cite{cho2015,sigmund2011}. The true diffusion coefficient $D_t$ is typically restricted in tumors due to high cellularity, which impedes free water diffusion~\cite{cho2015}. The pseudo-diffusion coefficient $D^{*}$ reflects microcirculatory blood flow and has been shown to differ between benign and malignant lesions, though it is the most challenging parameter to estimate due to its rapid signal decay, which is captured primarily within the first few low $b$-value acquisitions~\cite{lemke2011}.

Despite the rich physiological information encoded in IVIM parameters, their clinical adoption has been limited by the inherent instability of the biexponential fitting problem. The conventional approach voxel-wise non-linear least squares (NLLS) regression suffers from several critical limitations. First, the biexponential model is ill-conditioned: fitting three unknown parameters from a small number of noisy measurements (typically 6--12 $b$-values) leads to high estimation variance, particularly for the pseudo-diffusion coefficient $D^{*}$, whose exponential component decays rapidly and contributes meaningfully to the signal only at the lowest $b$-values~\cite{bertleff2017}. Second, NLLS operates independently on each voxel without leveraging spatial context, neighboring voxels within the same tissue type are expected to share similar parameter values, yet this prior information is entirely discarded. Third, at low signal-to-noise ratios (SNR), NLLS frequently fails to converge or converges to physically implausible parameter values, especially when the pseudo-diffusion coefficient upper bound is left unconstrained.

Recent advances in deep learning have demonstrated that neural networks can learn robust inverse mappings from noisy signals to underlying physical parameters, often outperforming classical optimization methods in ill-conditioned settings~\cite{barbieri2020,kaandorp2021}. Physics-informed neural networks, in particular, incorporate known physical laws directly into the network architecture or loss function, constraining the solution space to physically plausible regions and improving generalization to unseen noise conditions~\cite{raissi2019}. This paradigm is especially well-suited to the IVIM problem, where the exact forward model Equation~(\ref{eq:ivim}) is known analytically.

While several learning-based approaches to IVIM parameter estimation have been proposed~\cite{barbieri2020,kaandorp2021}, a systematic study of their robustness across a wide spectrum of clinically relevant noise levels, particularly in the context of breast cancer imaging where high-quality data may not always be available---remains lacking. Furthermore, most existing deep learning methods for IVIM operate on a voxel-by-voxel basis, ignoring the spatial structure of the parameter maps, which limits their denoising capability.

In this work, we present a two-stage physics-informed deep learning pipeline for noise-robust IVIM parameter estimation in breast DW-MRI. Our approach combines: (1) a Physics-Informed Autoencoder (PIA) that maps noisy DW-MRI signals to biophysical parameters using the known IVIM equation as its decoder, and (2) a U-Net-based spatial refiner that exploits spatial coherence in the parameter maps through tumor-weighted supervised denoising. We evaluate our pipeline against the conventional NLLS baseline across eight noise levels spanning SNR~$= 4$ to SNR~$= 100$, using 1,000 simulated breast phantom images from the VICTRE (Virtual Imaging Clinical Trials for Regulatory Evaluation) platform~\cite{badano2018}. Our key contributions are:

\begin{itemize}
    \item A self-supervised physics-informed autoencoder architecture (964,835 parameters) that constrains parameter estimates to biologically feasible ranges through a mean$\pm$delta$\cdot$tanh($\cdot$) formulation, eliminating the need for parameter labels during Stage~1 training.

    \item A spatially-aware CNN-PIA variant (189,443 parameters) that extends the voxel-wise MLP encoder with convolutional feature extraction, enabling spatial context exploitation at the Stage~1 level.

    \item A two-stage pipeline culminating in a U-Net spatial refiner (3,494,275 parameters) with residual learning and tumor-weighted supervision, achieving a total rRMSE of 0.12 at high SNR and 0.54 at extreme noise (SNR~$= 4$).

    \item A comprehensive noise-robustness evaluation on 400 held-out test patients demonstrating that the proposed CNN-PIA + Refiner pipeline maintains stable performance across all tested noise levels (SNR 4--100), whereas conventional NLLS diverges catastrophically at moderate-to-high noise ($\sigma \geq 0.10$), with total rRMSE exceeding 9.4 at SNR~$= 10$ compared to 0.28 for our method, a 33$\times$ improvement.
\end{itemize}

\begin{thebibliography}{12}

\bibitem{sung2021}
Sung, H., Ferlay, J., Siegel, R.\,L., et al. (2021).
Global Cancer Statistics 2020: GLOBOCAN Estimates of Incidence and Mortality Worldwide for 36 Cancers in 185 Countries.
\textit{CA: A Cancer Journal for Clinicians}, 71(3), 209--249.

\bibitem{lebihan1986}
Le Bihan, D., Breton, E., Lallemand, D., et al. (1986).
MR imaging of intravoxel incoherent motions: Application to diffusion and perfusion in neurologic disorders.
\textit{Radiology}, 161(2), 401--407.

\bibitem{lebihan1988}
Le Bihan, D., Breton, E., Lallemand, D., et al. (1988).
Separation of diffusion and perfusion in intravoxel incoherent motion MR imaging.
\textit{Radiology}, 168(2), 497--505.

\bibitem{cho2015}
Cho, G.\,Y., Moy, L., Kim, S.\,G., et al. (2015).
Evaluation of breast cancer using intravoxel incoherent motion (IVIM) histogram analysis: Comparison with malignant status, histological subtype, and molecular prognostic factors.
\textit{European Radiology}, 26(8), 2547--2558.

\bibitem{sigmund2011}
Sigmund, E.\,E., Cho, G.\,Y., Kim, S., et al. (2011).
Intravoxel incoherent motion imaging of tumor microenvironment in locally advanced breast cancer.
\textit{Magnetic Resonance in Medicine}, 65(5), 1437--1447.

\bibitem{lemke2011}
Lemke, A., Stieltjes, B., Schad, L.\,R., \& Laun, F.\,B. (2011).
Toward an optimal distribution of b values for intravoxel incoherent motion imaging.
\textit{Magnetic Resonance Imaging}, 29(6), 766--776.

\bibitem{bertleff2017}
Bertleff, M., Domsch, S., Weing\"{a}rtner, S., et al. (2017).
Diffusion parameter mapping with the combined intravoxel incoherent motion and kurtosis model using artificial neural networks at 3\,T.
\textit{NMR in Biomedicine}, 30(12), e3833.

\bibitem{barbieri2020}
Barbieri, S., Gurney-Champion, O.\,J., Klaassen, R., \& Thoeny, H.\,C. (2020).
Deep learning how to fit an intravoxel incoherent motion model to diffusion-weighted MRI.
\textit{Magnetic Resonance in Medicine}, 83(1), 312--321.

\bibitem{kaandorp2021}
Kaandorp, M.\,P.\,T., Barbieri, S., et al. (2021).
Improved unsupervised physics-informed deep learning for intravoxel incoherent motion modeling and evaluation in pancreatic cancer patients.
\textit{Magnetic Resonance in Medicine}, 86(4), 2250--2265.

\bibitem{raissi2019}
Raissi, M., Perdikaris, P., \& Karniadakis, G.\,E. (2019).
Physics-informed neural networks: A deep learning framework for solving forward and inverse problems involving nonlinear partial differential equations.
\textit{Journal of Computational Physics}, 378, 686--707.

\bibitem{badano2018}
Badano, A., Graff, C.\,G., Badal, A., et al. (2018).
Evaluation of digital breast tomosynthesis as replacement of full-field digital mammography using an in silico imaging trial.
\textit{JAMA Network Open}, 1(7), e185474.

\bibitem{gudbjartsson1995}
Gudbjartsson, H., \& Patz, S. (1995).
The Rician distribution of noisy MRI data.
\textit{Magnetic Resonance in Medicine}, 34(6), 910--914.

\end{thebibliography}

\end{document}
